% Authoritative LaTeX manuscript; Markdown is a generated compatibility projection.

% <!-- part-foreshadows: conversions-compound -->

\magevalue{chapter:software-factories} provided two kinds of evidence. The DocAble case followed one system over time, showing how engineering problems prompted changes to its models and controls. The company accounts described related mechanisms built independently under different constraints. Neither establishes a universal law. Together, they help us identify the engineering choices and tradeoffs that MAGE must address.

This chapter develops that account. We begin with production, lifecycle cost, and the environment inherited by successive engineering episodes. Four theoretical propositions then explain how representation leverage, independent enforcement, and engineering capital can shape the use of additional agentic capacity. Scope conditions identify where those relationships should matter, and a compact research agenda asks how to test them.

% <!-- principlebox -->

% <!-- box-family: canonical -->

\begin{quote}
\mbox{}%
\subsection{Key Ideas}\label{heading:the-synthesis}

The encoding–realization decomposition separates stating the intended obligations from fulfilling them. Engineering, however, does not consist of one attempt. Real systems are changed repeatedly, failures have consequences, and each episode can change the environment inherited by the next.

The question is how engineers make that inheritance useful: which knowledge to represent, which obligations to enforce, and which decisions should remain with a reasoner. The objective is durable progress at a justified cost, rather than maximum machinery or maximum implementation speed.
\end{quote}
